\section{Conclusion}
\label{sec:conclusion}

We studied the problem of incentivizing resource sharing between competing agents, without money, in settings where valuations are uncertain, 
revealed only through interaction, and sharing is irreversible. We formalized a novel model for incremental resource sharing, 
where a central planner proposes resource exchanges over multiple rounds. Our results show that carefully designed multi-round exchange 
algorithms can substantially outperform static policies by learning agents' valuations over multiple rounds. 

We characterized 
such optimal policies via a reduction to an optimal stopping problem on the lattice $\NN^2_{\geq 0}$. 
We then considered the problem of designing exchange algorithms when participants can leave early, and provided examples 
demonstrating why it is fundamentally more challenging.
These examples motivated necessary model conditions for which
optimal stopping and martingale methods can be used to characterize an optimal exchange algorithm.
% At the same time, we 
% identified fundamental limitations 
% when agents can strategically withdraw mid-policy, \kkcomment{expand}
Two open questions arise from our work. 
The first is whether or not an IC optimal algorithm is always guaranteed to exist, beyond the sufficient condition
we provide.
% While we have provided necessary conditions for optimal stopping and martingale methods
% to yield an IC algorithm, 
The second is how can our model and algorithms be extended to accommodate more than two competitors when there are complex interactions
between participants' utility functions, and thus collaboration incentives.

% Two open questions arising from our work is the  design of  \kkcomment{?} IC policies, \kkcomment{expand}, and the extension of our model 
% and methods to multiple agents.

% \section{IC outline}

% \begin{enumerate}
%     \item Section 4 outline
%     \item 
%     Say that the idea with optimal stopping is that once you reach a state, it doesn't matter how you got there, just where you are,
%     aka its kind of markovian. In these settings you can proceed greedily by pasting policies together.
%     However, this is not the case with IC policies because 
    
%     \item Give the 1/2 D example where IC pasting fails.

%     \item Argue that intuitively this represents a misalignment between welare and utility optimization.

%     \item In order to generalize the previous optimal stopping time algorithm approach to search for an IC policy, we need to make 
%     an assumption that IC policies can be combined in a gredy way. 

%     \item Assumption block

%     \item Under this assumption we can give a dp equation analog. First introduce new ess sup, State the proposition, say (-C) is a lower bound.
%     \item Proof sketch.

%     \item Explain that under this assumption we can contruct a policy that matches the performance of an optimal IC
    
%     \item Explain idea of theory, if we are allowed to stop vs not allowed to. Give theorem

%     \item Add remark that when this policy is IC, it is IC optimal.
    
%     \item Give proof citing proposition
    
%     \item \mycomment{NOte to self, revisit the stick note book marked, non trivially IC policy, see if it satisfies pasting 
%     assumptions? make sure it matches our policy?}

%     \item Add note that this really isn't practical to run. On going line of work to see if there are assumptions or special 
%     classes of welfare functions for which pasting is satisfied and its easier to run the alg

%     \item add corollary after main theorem saying if construction is IC then it is IC optimal
% \end{enumerate}  

% \begin{enumerate}
%     \item \mycomment{In the intro and techniques section probably needs a brief summary of what 
%     \citep{mandelbaum1981optimal} does so that the connection is not as abrupt}

%     \item 
%     say we're gonna start by reframing the problem. Say once an exchange of $(i,j)$ has occured, a know the value of the 
%     first i points and b knows the value of the first j points. To represent the information available after an exchange,
%     introduce sigma algebras for each of the states and define the family. And makke the observation that they are nested.
%     say intuitively this represents that you know more after a superset of the data has been exchanged. Initially, we know nothing
%     so $\Fcal_0=\cbr{\varnothing,\Omega}$

%     \item
%     % Say that given any algorithm, the corresponding sequence of states generated are a 
%     Say that corresponding to any alg there is the sequence (replace the $s_t$s with sigmas)
%     $\sigma_0,\ldots,\sigma_\tau$ which are a sequence of $\Scal$ 
%     valued random variables. More over, it is the case that $\sigma_{t+1}$ is ... measurable. This reflect that the choice 
%     the next exchange is only based on values observed from the previous exchange. 

%     \item Short proposition showing correspondence between algorithms and policies? 
%     \mycomment{feels unnecessary to make into prop but maybe reviewers would like to see more props?}

%     \item 
%     Give explanation saying that intuitively finding the optimal algorithm reduces to traversing 
%     the grid moving up/right at each step and determining when to stop. Therefore, we can use the 
%     language of (insert stopping paper) to formalize and analyze this problem.
%     The snell env is a classic concept from optimal stopping times 
%     We start with the snell env for partially ordered sets

%     \item Define snell envelope $X_s$ for $s\in\Scal$ (citing paper). Explain what it means 
%     intuitively, saying $\Pi_s$ is algorithms starting at state $s$.

%     \item Say in optimal stopping time theory problems, because $X_s$ is difficult to compute direclty,
%     you often prove a dynamic programming / recursive type of result. (cite paper) gives a formula
%     for paritially ordered sets, specializing to 0...m x 0...n gives (insert result for our
%     $N^2$ case). Then give proposition saying our optimal stopping algorithm?
%     % \item 
%     Have an algorithm block? Write it in algorithm form or policy form?
%     maybe instead of alg block, just have a theorem saying the specific algorithm is optimal 

%     \item 
%     Add remarks about how practically you would use techniques like MC methods to evaluate things 

%     \item Add note about how based on choice of welfare function you can make this indivisually
%     rational
%     \mycomment{add somewhere in assumptions that utility is 0 with no data exchanged.
%     without this IR being defined to be >0 doesn't make sense?
%     I think you only need to require $U^a_0=U^b_0=0$, don't worry about other states}
% \end{enumerate}  


% \subsection{An Optimal Algorithm when Participants cannot opt out}
% \paragraph{Reducting to an optimal stopping problem over N2 to solve}



% \begin{enumerate}
    
%     \item
%     Say how a algorithm can be viewed as a sequence of stopping points with respect to the 
%     sigma algebra generated by the values of data exchanged so far. Say that the stop decision
%     can be viewed as a stopping time w.r.t. the $\Gcal_t$.

%     \item
%     Define rewards for stopping as the welfare of the utilities.

%     \item
%     Therefore, finding the optimal data sharing algorithm reduces to finding the optimal stopping time
%     algorithm.

%     \item
%     Write out the algorithm to clairfy the interaction procedure

%     \item
%     Write a theorem saying that this is optimal. Proof involves using that rewards are bounded
%     and the stopping time is a.s. finite thus we can apply theorems 1 and 2 from the paper. 

%     \citet{mandelbaum1981optimal}

% \end{enumerate}  

% \paragraph{IR guarantees}

% \begin{enumerate}
%     \item Show that we automatically have ex ante IR for a sutible welfare function choice
%     under the optimal algorithm

%     \item 
%     Say we can also give a high probability ex post IR guarantee assuming we use the optimal
%     algorithm

%     \item Write and prove a proposition about this 
% \end{enumerate}  


% \section{Optimal algorithm when agents can opt}

% \begin{enumerate}
%     \item The model above assumes that agents either commit to following the data
%     sharing algorithm or do not, i.e. once they buy in they can't leave early.

%     \item But what happens if we assume agents can leave early. You might choose this model
%     for a real world setting if it is difficult to enforce agents to continue participating
%     or avoid agents who don't want to continue from submitting bad data
% \end{enumerate}  

% \subsection{Model changes when agents can opt out}

% \begin{enumerate}
%     \item Before agents made a single decision, whether or not to participate.
%     Now at each round they make a decision of whether or not to continue.
%     If either agent or the mechanism stops the sharing its over.
%     (Someone might ask what happens if both agents want to continue sharing but
%     the mechanism says to stop? In this case you can update the welfare function
%     so this won't happen?)

%     \item Define agent strategies as sequence of decisions

%     \item Remark that each agent's strategy corresponds to a stopping time
    
%     \item Use this to write out the expected welfare under a strategy profile in 
%     terms of a policy and agent strategy stopping times

%     \item Define a NE of policies under an algorithm

%     \item Define the set of IC policies using the min trick

%     \item State the problem of trying to find the optimal policy in this set,
%     i.e. the optimal IC policy.

% \end{enumerate}  

% \subsection{An Optimal Algorithm when Participants can opt out}


% \subsection{The opt IC algorithm}

% \begin{enumerate}
%     \item TBD
% \end{enumerate}  



% \section{Impossibility Results}

% \begin{enumerate}
%     \item 
%     Thus far we've assumed that agents report their observed values at each round
%     truthfully. The following result shows that if agents don't report truthfully
%     no algorithm can obtain the optimal expected welfare.

%     \item State proposition and give proof using two-flip.
    
% \end{enumerate}  


% \section{Running questions}

% \begin{enumerate}
%     \item I feel like there is a blind spot that I have for prior work in 
%     exchange of goods.

%     \item Should we consider other definitions of utility? It would make computations more annoying
%     but the results more general. (YES)

%     \item In the valuation functions example do I need to make sure everything is bounded? Technically to apply our
%     methods this is necessary but it would make the example cleaner if I didn't need to say $D_a,D_b$ are bounded and I
%     could write $V^a_0=V^b_0=-\infty$.

%     \item I don't even think independence of the value sequences are used anywhere?

% \end{enumerate}  

% \subsection{Small questions}

% \begin{enumerate}
%     \item Maybe give more examples of how value of data may not be known util it is used?

%     \item maybe show the calculations more in ex 1?

%     \item the general utility functions footnote?
    
%     \item should we specify probability space or require that the valuation functions are 
%     markov chains?

%     \item In the sharing alg def how much should we allude to the stopping time notation

%     \item Intellectual / property rights as another example of goods?

%     \item \mycomment{Do we need utilities to start at 0?}

%     \item
%     \mycomment{might definition to satisfy the property: to algs are = iff their dist over states is the same}

%     \item technically the paper gives a vague def for snell env as supermartingale ... and shows its equiv to 
%     what we say is the def

%     \item Weird that exchange policies is $\Pi^{\text{ex}}$ but IC policies, which is a subset, is $\Pi^{\text{IC}}$

%     \item fix stacking subscripts in pasting examples

%     % \item Is it ok that I chopen part of thm 1 from their paper?


% \end{enumerate}  

% \section{TODO}

% \begin{enumerate}
%     \item \mycomment{add assumption somewhere that tie breaking always chooses smaller or lexigraphically
%     smaller choice so that policies are consistent}
%     \item Resolve why the two flip calculation doesnt match the optimal policy calculation
%     \item cut down the descriptive text in section 2 
%     \item In existing results prove conditional dominated convergence,
%     proof is online
%     \item
%     \mycomment{add note about we must assume truthful reporting and its impossible to guarantee
%     this type of algorithm does better than doing nothing.
%     when agents are not truthful, welfare can be arbitrarily worse under the optimal algorithm } ?

%     \item Probably worth mentioning the connection of IC policies to solving for NE
%     \item the tildes and subscripts are hard to read, ex thm 3.5, fix this 
%     \item make sure not to take credit from paper in thm 3.5
%     \item I think the $\Ucal$ condition is probably redundant but harmless, maybe take this out later

%     \item I added the $\alpha$ policy notation in the truthfulness proof, make sure this is introduced somewhere

%     \item em dash creep

%     \item What should we call the sequence of exit strategies?

%     \item
%     \mycomment{In 4.2 to make more pasting more comparable, maybe make the 1d stopping time version not random and have it be stopping times which stop no 
%     earlier than fixed $t$}

%      \item \mycomment{we say that when random variables have finite support there are finitely many policies?
    %     is this only true assuming there is no external randomness to a policy?}
% \end{enumerate}  

% \section{To beef up submission}

% \begin{enumerate}
%     \item 
%     Add an equivalence proof of two definitions of algorithms

%     \item
%     Add the unbounded difference example

%     \item Work through the double flip example for the non-IC setting. write example followed 
%     by a proof block.

%     \item The IC alg might be more paletable if we can remove the $\Ucal$ part
    
%     \item 
%     Work through 4 flips example for the IC setting?

%     \item
%     Does pasting hold for reasonable welfare functions like min?


% \end{enumerate}  